\chapter{Závěr}

Cílem této diplomové práce bylo navrhnout a implementovat interaktivní aplikaci, která usnadní výuku a pochopení vybraných problémů z oblasti formální verifikace systémů. 
Konkrétní zaměření práce směřovalo na kalkulus komunikujících systémů, strukturální operační sémantiku a ohodnocené přechodové systémy. 
Tento cíl se podařilo v plném rozsahu naplnit vytvořením moderní a uživatelsky přívětivé webové aplikace.

V rámci teoretické a analytické části byly nejprve nastudovány principy reaktivních systémů a identifikovány nedostatky stávajících nástrojů v kontextu výuky začátků. 
Na základě těchto poznatků byla navržena aplikace, která uživatelům poskytuje okamžitou vizuální zpětnou vazbu a~provází je tak procesem formální modelace krok za krokem. 

Z technologického hlediska byla aplikace postavena na moderních tenchologiích, jako je React.js, TypeScript a Vite. 
Systém karet a modulární oddělení prezentační vrstvy od výpočetního jádra navíc zajišťují přehlednost, stabilitu a jednoduchou rozšiřitelnost celého řešení, což si myslím, že~je~velké plus této aplikace. 
Při vývoji aplikace jsem hlouběji pochopil problematiku komunikujících systémů, naučil jsem se ale také pracovat s několika moderními knihovnami, jako je Peggy.js a~Cytoscape.js. 
Jak již bylo nastíněno, užitečným rozšířením aplikace by mohla být implementace modulů pro ověřování silné a slabé ekvivalence chování (bisimulace) či integrace nástrojů pro model checking s~využitím Hennessy-Milnerovy logiky.


\endinput